\PassOptionsToPackage{table}{xcolor}
\documentclass{article}
\usepackage{iclr2027_conference,times}
\usepackage{hyperref}
\usepackage{url}
\usepackage{graphicx}
\usepackage{booktabs}
\usepackage{amsmath}
\usepackage{xcolor}
\definecolor{linkblue}{HTML}{1F5A94}
\hypersetup{
  colorlinks=true,
  linkcolor=linkblue,
  citecolor=linkblue,
  urlcolor=linkblue
}
\usepackage{amsfonts}
\usepackage{amssymb}
\usepackage{amsthm}
\usepackage{enumitem}
\usepackage{pifont}
\DeclareMathSizes{9}{10}{7}{5}

% Theorem environments and shared math notation.
\newtheorem{theorem}{Theorem}
\newtheorem{assumption}{Assumption}
\newtheorem{proposition}{Proposition}
\theoremstyle{definition}
\newtheorem{definition}{Definition}
\theoremstyle{plain}
\newcommand{\E}{\mathbb{E}}
\newcommand{\Var}{\mathrm{Var}}
\newcommand{\norm}[1]{\lVert #1 \rVert}

% Tighten float/caption spacing to fit the 9-page limit (formatting only).
\setlength{\textfloatsep}{8pt plus 2pt minus 2pt}
\setlength{\floatsep}{8pt plus 2pt minus 2pt}
\setlength{\intextsep}{8pt plus 2pt minus 2pt}
\setlength{\abovecaptionskip}{4pt}
\setlength{\belowcaptionskip}{2pt}

% Compact captions: keep plots readable while avoiding body-sized caption blocks.
\makeatletter
\long\def\@makecaption#1#2{%
  \vskip\abovecaptionskip
  {\footnotesize
    \sbox\@tempboxa{\textbf{#1:} #2}%
    \ifdim \wd\@tempboxa >\hsize
      \textbf{#1:} #2\par
    \else
      \global\@minipagefalse
      \hb@xt@\hsize{\hfil\box\@tempboxa\hfil}%
    \fi}
  \vskip\belowcaptionskip}
\makeatother

% -----------------------------------------------------------------------
% Working title — change before submission
% -----------------------------------------------------------------------
\title{Weights Have Traits: Estimating Fine-Tuning Lineage in LLMs}

% Authors hidden for double-blind review.
% \iclrfinalcopy  % Uncomment for camera-ready only.

\newcommand{\fix}{\marginpar{FIX}}
\newcommand{\new}{\marginpar{NEW}}

% Placeholder macro for unwritten sections — shows in red, visible in draft
\newcommand{\TODO}[1]{\textcolor{red}{\textbf{[TODO: #1]}}}

% Shorthand for figure path prefix
\newcommand{\figpath}{figures}

\begin{document}
\maketitle

% -----------------------------------------------------------------------
\begin{abstract}
While over a million deep learning models are hosted online, most are fine-tuned descendants from a small set of ancestors.  The lineage of these models is only sometimes documented fully, even though provenance, auditing, and safe deployment all depend on proper documentation. In response to this issue, we hypothesize that fine-tuning acts as an evolutionary process that leaves signal in weight space, and show that a model \textit{phylogeny} (evolutionary tree) can be estimated from model weights alone and scored against a known ground-truth phylogeny, recovering all true clades (branches) in 85--100\% of runs.  We treat individual weight tensors as gene-like estimators of lineage, revealing where evolutionary signal is concentrated and allowing each tensor to produce its own phylogeny.  We are able to show that we can accurately estimate the phylogeny using only a percent of weights (1-3\% in our experiments), suggesting that lineage can be recovered from a small, architecture-specific weight subset. Moreover, we show that this phylogeny estimation is associated with known mathematical concepts from phylogenetics: the Atteson margin and four-point additivity. Beyond phylogeny estimation, we report two additional empirical findings. First, weight-space proximity predicts behavioral similarity on several held-out tasks involving truthfulness, knowledge, reasoning, and common sense, though the strength and statistical support depend on the probe. Second, weight-based phylogeny estimation can match or outperform behavior-based estimation, particularly after low-rank updates. We finally compare our weight-based method to a state-of-the-art behavior-based method on both our controlled experiments and a real, albeit small, modern model phylogeny. On the matched controlled trees, weights recover 100\% of clades versus PhyloLM's 35\% under rank-8 LoRA and 100\% versus 52\% under full fine-tuning, while both methods recover all three documented splits in the real Mistral family.

\end{abstract}

% -----------------------------------------------------------------------
\section{Introduction}
\label{sec:intro}




While there are over one million models available on HuggingFace for deep learning and a handful of frontier
labs constantly releasing updated versions, most models are
fine-tuned descendants of a few ancestors.  Very few of those models transparently document their full model lineage \citep{horwitz2025chartatlasworldsmodels}. Benchmarks can reveal how a model performs on a task, but not how a model originated, how it differs from its sibling models, or how it may
fail. In contrast, training lineage can reveal much about model provenance, auditing, and safety~\citep{behavioral2026provenance, nikolic2026model, tensorguard2025}.

We, like other prior works~\citep{horwitz2025mother, yu2025neural, wu2026llmdna,gallagher2026}, observe that fine-tuning is a path-dependent, branching process. Each
child model inherits its parent's  updated weights, nudged towards a certain loss and/or data set. This is
precisely an evolutionary process, as described in \cite{yax2025phylolm}, and moreover, the evolutionary path leaves signal in weight space. We treat a line or genealogy of fine-tuned models
as a phylogeny (i.e. an evolutionary tree~\cite{saitou1987neighbor}) and ask whether this phylogeny can be \emph{estimated} from the model weights alone.

As motivation, consider a firm that forks a base model into finance, IT, and fraud variants, each fine-tuned until
dozens of versions circulate with no recorded lineage (Fig.~\ref{fig:overview}A). A practitioner needs
to know which to build on, where two versions diverged, and where each fails. Phylogenetic algorithms
were built for exactly this inference in biology~\citep{saitou1987neighbor, robinson1981comparison},
and we adapt them to neural net model weights. We estimate a phylogeny from model weights and score it against a \emph{known}
ground-truth phylogeny (Fig.~\ref{fig:overview}A), link the changes in weights to changes in behavior (Fig.~\ref{fig:overview}B), and ask whether diagnostic statistics indicate when phylogeny estimation will succeed even without access to the true tree (Fig.~\ref{fig:overview}C).

\begin{figure}[t]
  \begin{center}
    \includegraphics[width=\linewidth]{\figpath/fig1_paper_overview.pdf}
  \end{center}
  \caption{Overview of our three questions: Can model weights recover known training relationships? Do nearby models behave alike? Do tree-like distances signal reliable recovery without knowing the true tree?}
  \label{fig:overview}
\end{figure}



Our primary contribution is that analyzing weight changes between fine-tunings allows accurate phylogeny estimation under the assumption of latent ancestors. We use a fitted Atteson diagnostic and label-free four-point additivity as empirical diagnostics to characterize when phylogeny estimation succeeds. As a result of these weight-based estimates, we can show two behavior-related consequences. In summary we show:
\begin{itemize}[leftmargin=1.3em,itemsep=1pt,topsep=2pt,parsep=0pt]
  \item \textbf{Phylogeny estimation from leaf model weights against ground truth.} Guided by the geometry of
    fine-tuning updates, our distance metric and summary-statistic choices estimate the model phylogeny from
    weights alone. Our controlled experiments recover \textit{all} true clades (branches) in $85$--$100\%$ of runs
    ($92$--$100\%$ clade recovery). We evaluate against known ground-truth phylogenies and show that a subset containing about 1-3\% of the scalar weights can accurately estimate the known phylogenies. Treating weight tensors as gene-like estimators also yields a separate phylogeny from each tensor and signals where lineage signal is prominent. We show that our weight-based method accurately recovers all the splits in a small but real Mistral phylogeny.

\item \textbf{Scenarios where weight-based estimation outperforms behavior-based estimation.}
    On the same 46 controlled trees, weights recover 100\% of clades versus 35\% for
PhyloLM~\citep{yax2025phylolm} under rank-8 LoRA, and 100\% versus 52\% under full fine-tuning.
    On a documented Mistral family with large RLHF/DPO forks, both methods
    recover all three true splits. Thus behavior-only estimation can succeed after behaviorally large
    updates, while weight access remains substantially more informative for smaller updates and
    requires no output sampling.
  \item \textbf{Evidence of similar weights signaling similar behavior.} Changes in weights can be used to predict similarity in model behavior.  In particular, we show that models that are closer in weight-space behave more similarly in behavior-space on \textit{held-out} tasks such as truthfulness, reasoning, knowledge, and common sense.

\end{itemize}
% Submission TODO: add the anonymous code-and-artifact URL here once the blind snapshot exists.
Our main controlled topology results cover Flan-T5-base (250M) and Llama-3.2-1B, with both full
fine-tuning and variants of LoRA.






% -----------------------------------------------------------------------
\section{Related Work}
\label{sec:related}
% ----------------------------------------------------------------------

Deep learning lineage studies take three broad forms: \emph{lineage from model weights}, \emph{lineage from behavior}, and \emph{connections between weight and behavior changes}. Table~\ref{tab:related} compares some of the closest methods. These studies are driven by applications in provenance and IP
\citep{horwitz2025mother, zeng2024huref, nikolic2026model}, model poisoning and misalignment
\citep{cloud2025subliminal}, trait inheritance \citep{cloud2025subliminal, agents2026same}, and data
auditing \citep{zhu2025independence}. 

\textbf{Model weights.} Several works trace weight changes from fine-tuning or reinforcement learning with human feedback (RLHF) back to parent models. MoTHer
\citep{horwitz2025mother} recovers directed model trees from weight similarity via a directed minimum
spanning tree and evaluates against controlled, known model hierarchies; \citet{yu2025neural} recover parent--child and grandparent--child lineage with both a
learning-free fine-tuning approximation and a learned detector ($>99\%$ accuracy). Beyond direct
lineage, spectral weight ``fingerprints'' identify common ancestry \citep{ghostspec2025}, and Spectral
DeTuning \citep{horwitz2024spectral} recovers base weights from shared LoRA variants. Our setting
differs from MoTHer's observed-node model tree in that we observe only terminal nodes and estimate a
leaf phylogeny with latent rather than explicit internal ancestors.  This allows for estimation of phylogenies as opposed to directed spanning trees, both of which can have practical applications in the field.

Other works examine \emph{why} lineage is detectable in models. Models fine-tuned from a shared initialization stay
within a single low-error basin and are robust to SGD noise
\citep{neyshabur2020,wortsman2022model,frankle2020linear}, so pairwise weight distance can approximate genealogical distance. Additionally, task vectors \citep{ilharco2023editing} are near-orthogonal across tasks, so
distinct lineages contribute independent directions. LoRA \citep{hu2022lora} and intrinsic-dimensionality
results \citep{aghajanyan2021intrinsic} help explain why small-rank updates suffice and why rank matters.



\textbf{Model behavior.} An advantage of behavior only studies is they need only API access or benchmark outputs
\citep{yax2025phylolm, wu2026llmdna} as opposed to full model access. For example, PhyloLM \citep{yax2025phylolm} builds LLM phylogenies from output
similarity and predicts benchmark performance; LLM DNA \citep{wu2026llmdna} defines a functional
``DNA'' representation carrying lineage. Both methods reconstruct phylogenies over hundreds (or many more) of models. Other
output-only methods test pairwise provenance from next-token distributions \citep{nikolic2026model} or
output independence \citep{zhu2025independence}. These behavior-only methods are advantageous as they do not require shared
architecture and can compare models across families and sizes, an advantage weight-based methods currently lack.
These studies, however, are typically unsupervised and prompt-dependent, with no ground-truth phylogeny to score
against, making it difficult to quantify their success.

\textbf{Model weights $\leftrightarrow$ behavior.} A second group links weight changes to
behavior. \citet{cloud2025subliminal} show child models inherit ``subliminal'' traits even when every
reference to the trait is filtered from training text, and prove a parent's preference is preserved
under a large enough training signal given a shared initialization. REEF \citep{zhang2025reef} compares
internal representations (via centered kernel alignment CKA~\citep{kornblith2019similarity}) to decide lineage, and gradient-response fingerprints
\citep{tensorguard2025} classify families from how weights react to probes. These establish that
weight closeness are linked to behavioral similarity, but only pairwise, not as a full phylogeny.

We bridge weight- and behavior-based lineage, determine where the phylogeny can be estimated, and evaluate against ground-truth trees. Our findings are consistent with prior work showing that inherited behavioral signals can be preserved under sufficient training signal \citep{cloud2025subliminal}, as well as phylogenetic results showing that the underlying tree can be correctly reconstructed when certain conditions on the estimated distances are satisfied \citep{atteson1999performance}.

\begin{table}[t]
\caption{Comparison with the closest lineage-estimation methods.}
\label{tab:related}
\begin{center}
\small
\setlength{\tabcolsep}{6pt}
\begin{tabular}{lcccc}
\toprule
Method & Access & \shortstack{Latent\\ancestors} & \shortstack{Known-lineage\\eval} & \shortstack{Real\\checkpoints} \\
\midrule
MoTHer \citep{horwitz2025mother}
  & weights  & \ding{55}  & \checkmark & \checkmark \\
Neural Phylogeny \citep{yu2025neural}
  & weights  & \ding{55}   & \checkmark & \checkmark \\
PhyloLM \citep{yax2025phylolm}
  & behavior & \checkmark  & \ding{55}  & \checkmark \\
LLM DNA \citep{wu2026llmdna}
  & behavior & \checkmark  & \ding{55}  & \checkmark \\
\textbf{Ours}
  & \textbf{both} & \checkmark & \checkmark & \checkmark \\
\bottomrule
\end{tabular}
\end{center}
\end{table}



% -----------------------------------------------------------------------x
\section{Framework}
\label{sec:framework}
% -----------------------------------------------------------------------

We first define our terms, notation, and statistical framework. A \emph{phylogeny} $\mathcal{P}$ is a tree
whose leaves are the observed models and whose internal nodes are latent common ancestors; a
\emph{clade} is the set of leaves descending from one such ancestor (one branch of the tree). We analyze
single-root model families but observe only their terminal models. Thus,  we estimate the unrooted leaf topology with internal ancestors latent. We form a pairwise distance matrix and apply neighbor-joining (NJ)
\citep{saitou1987neighbor} to estimate the phylogeny.

\textbf{Phylogeny estimation.}\label{sec:framework:phylo}
Assume we have a model $g(\theta)$ where $\theta$ is composed of multiple layers $\theta_\ell$. To estimate $\widehat{\mathcal{P}}$: (1) choose a set of layers $\theta_\mathcal{L}$ and a distance metric $d$;
(2) for each layer $\ell$ compute $D_{i,j}^{(\ell)} = d(\theta^{(\ell)}_i, \theta^{(\ell)}_j)$ between models
$i,j$; (3) assemble $\widehat{\mathbf{D}}$ as a function of the $D_{i,j}^{(\ell)}$ (e.g.\ the across-layer
average); and (4) estimate $\widehat{\mathcal{P}} = f(\widehat{\mathbf{D}})$ with $f$ as the NJ method.

Step (3) lets different layers yield different estimates of $\widehat{\mathcal{P}}$, which we use to
assess each layer's lineage signal. We score an estimated phylogeny against the known truth with the
Robinson--Foulds (RF) distance \citep{robinson1981comparison}, the number of bipartitions on which the
two trees disagree, which decomposes into false positives (FP) and false negatives (FN). Because NJ
returns strictly binary trees while our ground-truth phylogenies contain \textit{polytomies} (nodes with
more than two children), NJ is forced to add resolving edges. As a result, RF is never zero even for a correct binary
reconstruction. These forced edges inflate FP without being genuine errors. We therefore treat FN, the
number of true clades the reconstruction misses, as a more representative metric. Additionally, we report clade
recovery $= 1 - \mathrm{FN}/|\text{true splits}|$ and polytomy-aware exact recovery (PAER), the fraction of
runs with $\mathrm{FN}=0$, alongside RF. We also report the standardized within-run branch correlation
$r_{\mathrm{branch}}$ between weight distance and a same-branch indicator. We analyze only unrooted
phylogenies with at least four leaves because three leaves admit only one unrooted topology and no
nontrivial internal split.


\textbf{Distance metric.}\label{sec:framework:metric}
We flatten each weight tensor to a vector and compare models with vector distances rather than matrix
norms, for portability across architectures, considering cosine, correlation, $\ell_1$, and $\ell_2$
(the last used in MoTHer~\citep{horwitz2025mother}). We adopt cosine distance ($1$ minus cosine
similarity) throughout, for three reasons. It is bounded and scale-invariant across layers, so no single
high-magnitude layer dominates the aggregate, unlike $\ell_1$ or $\ell_2$. It is the prevailing choice in
recent provenance and fingerprinting work \citep{zeng2024huref, behavioral2026provenance,
zhu2025independence}, keeping our distances comparable to that literature. Finally, it has a direct geometric
reading on merged checkpoints: descendants whose weight vectors are deflected in similar directions remain closer in angle. For LoRA-trained trees, we merge the ordered root-to-leaf adapter chain into each checkpoint before comparing weights.



\textbf{Per-tensor estimated phylogenies.}
We treat each weight tensor as a gene-like estimator rather than collapsing immediately to a single
distance. This is analogous to biological phylogenies that must cohere across genes encoding different
information~\citep{maddison1997gene}. Each analyzed tensor yields one cosine-distance matrix and one
NJ tree. Our primary topology estimate is the Billera--Holmes--Vogtmann (BHV) Fr\'echet mean of these
trees, i.e., the tree minimizing the sum of squared BHV geodesic distances
~\citep{billera2001geometry, willis2019confidence}. We include only tensors with nonzero between-model
cosine distance; an exactly unchanged LoRA tensor, or a changed tensor whose vectors differ only by
positive rescaling, has an all-zero cosine-distance matrix and does not define a tensor-level phylogeny.
The controlled analysis uses the original cosine-distance
branch-length scale; normalizing each active tensor by its mean pairwise distance gives a sensitivity
analysis. We retain the arithmetic mean of the active-tensor distance matrices for pairwise lineage,
four-point additivity, and the fitted Atteson diagnostic, all of which require a dissimilarity matrix.
This tensor-wise view both provides an aggregate topology and lets us rank where inherited signal is
carried (\S\ref{sec:results:weighttree}).


\textbf{When estimation succeeds.}\label{sec:framework:theory}
Model lineage can be probed at two levels: pairwise ordering (e.g.\ whether two models sit on the same
branch) or whole-tree recovery, the former being easier. For ordering, we pool all pairwise model
comparisons within each valid run and test with the Mann--Whitney $U$ statistic. This branch-level
diagnostic is simpler than either pairwise parent--child detection~\citep{yu2025neural}, MoTHer's
directed observed-node tree recovery~\citep{horwitz2025mother}, or our latent-ancestor leaf topology. Recovering the full
\textit{phylogeny} is harder but well studied in phylogenetics. The NJ method is exact when the distance matrix is
additive, a condition model-weight distances satisfy only approximately. Atteson's theorem makes the
tolerance precise.
\begin{theorem}[\citealp{atteson1999performance}]\label{thm:atteson}
Let $D_T$ be the additive leaf metric of an edge-weighted binary tree $T$ whose edges have strictly
positive lengths $\{L_s\}$, and let $D$ be an observed dissimilarity map. If
$\lVert D-D_T\rVert_\infty < \tfrac12\min_s L_s$, then neighbor-joining applied to $D$ returns $T$.
\end{theorem}
For the theorem's fixed true metric, the ratio
$m:=\min_s L_s/(2\lVert D-D_T\rVert_\infty)>1$ is a sufficient condition of exact recovery. It is
conservative, so NJ can recover the topology when $m\leq1$. In our controlled experiments the true topology is
known but its edge lengths are not. We fit nonnegative edge lengths on that topology to obtain
$D_{\mathrm{fit}}$ and report the analogous diagnostic
$\widetilde m=\min_s\widehat L_s/(2\lVert D-D_{\mathrm{fit}}\rVert_\infty)$. Because
$D_{\mathrm{fit}}$ is estimated from $D$ and some truth trees contain polytomies, $\widetilde m$ is an approximation rather than the true Atteson condition.  Because of this fitting approximation, we also
adopt a label-free proxy computed from $D$ alone, based on the classical four-point condition for tree
metrics~\citep{sempleSteel2003phylogenetics}.

\textit{Four-point additivity} ($A$). Non-additivity of a distance matrix is measured by the quartet
$\delta$-score of \citet{holland2002delta}. For each quartet of leaves the three pairwise-distance sums,
sorted $s_1\ge s_2\ge s_3$, give
\[
  \delta_q=\frac{s_1-s_2}{s_1-s_3}\in[0,1] ,
\]
which is $0$ exactly when the four-point condition $s_1=s_2$ holds \citep{sempleSteel2003phylogenetics}
(the additive, tree-metric case) and $1$ at the maximally non-treelike quartet (the two \emph{smallest}
sums equal, the largest strictly greater). We report its complement, averaged over the $Q$ quartets,
\[
  A=\frac{1}{Q}\sum_{q=1}^{Q}\bigl(1-\delta_q\bigr)
   =\frac{1}{Q}\sum_{q=1}^{Q}\frac{s_2^{(q)}-s_3^{(q)}}{s_1^{(q)}-s_3^{(q)}}\in[0,1] ,
\]
so that $A=1$ is perfectly additive and larger $A$ is associated with more accurate phylogeny estimation. Because $A$ requires no labels,
it serves as our proxy when the true tree is unknown. For evaluation against a known tree we
additionally report $A$ over only the \emph{informative} quartets (those the truth resolves into two
cherries), since quartets that fall inside a polytomy carry no split to recover and only dilute the
all-quartet mean. We relate $A$ and the fitted Atteson diagnostic to phylogeny estimation accuracy in \S\ref{sec:results}.


% -----------------------------------------------------------------------
\section{Experiments}
\label{sec:experiments}
% -----------------------------------------------------------------------
We structure our experiments around three questions: whether the model lineage or phylogeny can be
estimated from weight changes; whether those changes predict downstream held-out behavior; and
how behavior-based phylogeny differs from weight-based.

\paragraph{Controlled and real training trees.}
\label{sec:exp:trees}
% Base: Flan-T5-base (250M). Bernoulli branching, 4 tasks, 50 replications,
% 5--8 leaves, known topology by construction.
% 2000 steps, lr=3e-4, batch=8, grad_accum=4, bf16.
% Controlled design enables exact evaluation against the latent-ancestor leaf topology.

To create a controlled experiment, we generate stochastic model training trees beginning from a root model (process detailed in Appendix~\ref{app:details}), producing 50 trees with up to 10 leaf nodes. We randomly assign a task/dataset pair to each edge, and at each split fine-tune a child model from its parent on that pair. LoRA chains are sequentially merged so the primary LoRA leaves are standalone full-weight checkpoints; the separately marked all-projections ablation is analyzed in cumulative update space. We then estimate a phylogeny from the \textit{model leaves} and compare it to the known training tree, testing both direct lineage and phylogeny estimation.

We vary the root model, task types, and trained layers. Our primary evaluation uses an
encoder--decoder model (Flan-T5-base, 250M) and a decoder-only model (Llama-3.2-1B), fine-tuning on
Q\&A, classification, translation, and summarization. We also run auxiliary analyses on Pythia-160M and Llama-3.1-8B. For Flan-T5, we run full fine-tuning and four LoRA \citep{hu2022lora}
scopes (key only; query/key/value; query/key/value/output; all attention and feed-forward projections).
For Llama-3.2-1B, we run full fine-tuning and query/key/value LoRA at ranks 8 and 64. For each base
model, we use the same task instructions. We also
validate our method on three documented model lineages from HuggingFace. We highlight one in the
main text and report the other two in Appendix~\ref{app:organic-additional}.

\paragraph{Direct lineage.}
MoTHer~\citep{horwitz2025mother} assembles pairwise relatedness and direction scores into a directed
model tree, while Neural Phylogeny~\citep{yu2025neural} identifies parent--child relationships. Our
auxiliary ordering analysis instead works at the branch level among observed leaves. For each run we test whether same-branch leaf pairs are closer in weight space than cross-branch pairs with a one-sided Mann--Whitney $U$ test, summarize the effect with the rank-biserial correlation, and report the within-run branch correlation $r_{\mathrm{branch}}$ (Fisher $z$-averaged across runs)\footnote{For each run we compute $z=\operatorname{arctanh}(r)$, average on the $z$ scale, and back-transform with $r=\tanh(\bar z)$.} between weight distance and the same-branch indicator.

\paragraph{Phylogeny estimation.}
\label{sec:exp:weights}
Besides direct lineage, we test whether our method estimates the phylogeny. From \S\ref{sec:framework:theory}, accurate estimation should be possible when the distance matrix is sufficiently additive, for which the label-free four-point additivity $A$ is a practical, label-free proxy. For each tree we report RF, FN, clade recovery, and polytomy-aware exact recovery (PAER; $\mathrm{FN}=0$), relating them to the fitted Atteson diagnostic $\widetilde m$ and $A$ in \S\ref{sec:results}. The per-tensor trees also let us identify which tensors carry the most inherited signal, via the highest- and lowest-FN tensors per architecture.


\paragraph{Predicting held-out behavior.}
\label{sec:exp:behavioral}
We next estimate how strongly weight distance predicts behavioral similarity on probe tasks using the trained, controlled experiment trees. For the translation probe, we obtain one deterministic greedy response per prompt and model, embed each output with a small sentence-transformer (\texttt{all-MiniLM-L6-v2}, $384$-d), and average prompt-paired cosine similarity. We also evaluate HellaSwag, ARC-C, MMLU, and TruthfulQA on the same trees using the corresponding prompt-paired semantic endpoint. We then report the random-effects DerSimonian--Laird \citep{dersimonianlaird1986} pooled correlation $r_{\mathrm{DL}}$ between weight distance and output-embedding similarity, with both quantities $z$-scored within run (placing all models and architectures on one scale). Other generation details are reported in the appendix.




\paragraph{Comparison to behavior-only estimation.}
\label{sec:exp:blackbox}
Sometimes weights are unavailable and provenance must be inferred from outputs alone, like the method PhyloLM \citep{wu2026llmdna,yax2025phylolm}.  Because our documented HuggingFace family (\S\ref{sec:exp:trees}) has a known topology, we can score behavior-only estimation against its documented reference lineage. On that family we run the behavior-based PhyloLM method~\citep{yax2025phylolm}, using its released genes and genetic distance on the same checkpoints, and compare it against our weight-based estimation on the identical documented lineage (details in App.~\ref{app:blackbox}). Our controlled comparison runs PhyloLM on every tree for the Llama rank-8 LoRA, rank-64 LoRA, and full-FT conditions (50 reconstructions per condition). We compare on trees with at least one nontrivial unrooted truth split. The encoder--decoder Flan groups are outside PhyloLM's decoder-only scope.


% -----------------------------------------------------------------------
\section{Results}
\label{sec:results}
% -----------------------------------------------------------------------

%\subsection{Weights contain both direct ancestry and phylogenetic signal}

We examine the results of our controlled and real data experiments with our encoder-decoder and decoder models, LoRA trainings, and different task fine-tunings and data sets.


\paragraph{Direct lineage.}
We first determine how well direct lineage estimation (e.g., parent--child or within-branch) performs under different conditions. There are $n_{ord}=26$ eligible trees for this analysis that contain at least one pair of leaves from the same root branch and at least one pair from different root branches. Table~\ref{tab:variants} reports the Mann--Whitney rank-biserial effect and the within-run branch correlation $r_{\mathrm{branch}}$ between cosine distance and a same-branch indicator. Every mean rank-biserial effect is positive and every $r_{\mathrm{branch}}$ is negative. This implies that ordering in a tree has a strong signal in the weights.

\begin{table}[h]
\caption{Lineage and recovery by training regime (cosine; mean with SE). Topology recovery uses the
  BHV Fr\'echet mean of per-tensor NJ trees; pairwise lineage statistics use the mean active-tensor
  distance matrix. LoRA uses sequentially merged full weights and only tensors with nonzero
  between-model cosine distance;
  $\dagger$ marks the cumulative update-space all-projections ablation.
  PAER means $\mathrm{FN}=0$. Llama weight and PhyloLM values share the same $n_{\rm rec}=46$
  eligible trees; Flan is outside PhyloLM's scope. $^\star$Ceiling cells correspond to 117/117
  recovered truth splits and 46/46 PAER trees (two-sided exact 95\% binomial intervals
  96.9--100\% and 92.3--100\%, respectively).}
\label{tab:variants}
\begin{center}
\scriptsize
\setlength{\tabcolsep}{1.1pt}
\resizebox{\linewidth}{!}{%
\begin{tabular}{lrrrrrrr}
\toprule
Variant & $n_{\rm rec}/n_{\rm ord}$ & \shortstack{Rank-bis.\\(SE)} & \shortstack{$r_{\mathrm{branch}}$\\(SE)} & \shortstack{Weight clade\\(SE)} & \shortstack{PhyloLM clade\\(SE)} & \shortstack{Weight PAER\\(SE)} & \shortstack{Weight RF (SE);\\FN (SE)} \\
\midrule
\multicolumn{8}{l}{\textbf{Encoder--decoder: Flan-T5-base} ($d=768$)} \\
\multicolumn{8}{l}{\quad\textit{LoRA adapter scope}} \\
\quad LoRA: key only                    & 46/26 & 0.52 (0.09) & $-0.53$ (0.07) & \textbf{92\%} (3\%) & n/a & \textbf{85\%} (5\%) & 2.39 (0.25); 0.17 (0.06) \\
\quad LoRA: q/k/v                       & 46/26 & 0.56 (0.09) & $-0.62$ (0.07) & \textbf{99\%} (1\%) & n/a & \textbf{98\%} (2\%) & 2.09 (0.21); 0.02 (0.02) \\
\quad LoRA: q/k/v/o (all attention)     & 46/26 & 0.56 (0.09) & $-0.62$ (0.07) & \textbf{99\%} (1\%) & n/a & \textbf{98\%} (2\%) & 2.09 (0.21); 0.02 (0.02) \\
\quad LoRA: q/k/v/o + FFN projections$^\dagger$ & 46/26 & 1.00 (0.00) & $-0.91$ (0.02) & \textbf{96\%} (2\%) & n/a & \textbf{89\%} (5\%) & 2.26 (0.21); 0.11 (0.05) \\
\multicolumn{8}{l}{\quad\textit{Full fine-tuning}} \\
\quad Full fine-tuning                  & 46/26 & 0.69 (0.08) & $-0.71$ (0.06) & \textbf{100\%}$^\star$ & n/a & \textbf{100\%}$^\star$ & 2.04 (0.19); 0.00 (0.00) \\
\midrule
\multicolumn{8}{l}{\textbf{Decoder-only: Llama-3.2-1B} ($d=2048$)} \\
\multicolumn{8}{l}{\quad\textit{Controlled exp. trees}} \\
\quad LoRA qkv ($r{=}8$)                 & 46/26 & 0.63 (0.08) & $-0.64$ (0.07) & \textbf{100\%}$^\star$ & 35\% (5\%) & \textbf{100\%}$^\star$ & 2.04 (0.19); 0.00 (0.00) \\
\quad LoRA qkv ($r{=}64$)                & 46/26 & 0.65 (0.08) & $-0.67$ (0.07) & \textbf{100\%}$^\star$ & 24\% (5\%) & \textbf{100\%}$^\star$ & 2.04 (0.19); 0.00 (0.00) \\
\quad Full fine-tuning                   & 46/26 & 0.64 (0.08) & $-0.66$ (0.07) & \textbf{100\%}$^\star$ & 52\% (5\%) & \textbf{100\%}$^\star$ & 1.83 (0.17); 0.00 (0.00) \\
\bottomrule
\end{tabular}
}
\end{center}
\end{table}



\paragraph{Phylogeny estimation.}
\label{sec:results:weighttree}
In addition to ordering, the BHV Fr\'echet means estimate the phylogeny, as shown in Table~\ref{tab:variants}, with respect to clade recovery, polytomy-aware exact recovery (PAER; $\mathrm{FN}=0$), and RF and FN for all eligible phylogenies. In the controlled experiments, clade recovery ranges from 92--100\% and PAER from 85--100\%. The merged LoRA variants recover 92--99\% of clades for Flan and 100--100\% for Llama; full fine-tuning recovers 100\% and 100\%, respectively. On the same 46 trees with a nontrivial truth split, weights recover 100\% versus PhyloLM's 35\% under rank-8 LoRA, 100\% versus 24\% at rank 64, and 100\% versus 52\% under full fine-tuning. We see in Appendix~\ref{app:withintask} that phylogeny estimation from weights still holds when every node in a tree trains within a single task (clade recovery $59$--$87\%$); in that design, task-family identity provides no branch-specific cue, although path length can retain leaf-depth information.

Treating tensors as gene-like estimators reveals that a small attention-weight subset can estimate the phylogeny as well as the all-tensor aggregate in the full-fine-tuning regime. On the same 46 matched-topology Flan runs, the encoder self-attention key subset and the all-weight estimate both recover 100\% of clades with 100\% PAER (App.~\ref{app:layers}). This subset contains 7.08M of 247.58M scalar weights (2.86\%, approximately 3\%). The individual-tensor profile in App.~\ref{app:layertrace} is more diffuse: Flan key matrices and other tensors each miss about 4\% of true clades on average.

Finally, Figure~\ref{fig:coherence_recovery} shows simulation-level, tensor-subset, and individual-matrix recovery summaries from the same full-weight trees as Table~\ref{tab:variants}, together with the matched PhyloLM estimates for the three Llama conditions. Weight-based estimation of clade recovery ranges from 92--100\%, including 100\% for Flan full fine-tuning and 100\% for Llama full fine-tuning. Tensor-subset estimates recover 90--100\% of true clades. The matched PhyloLM estimates recover 24--52\% of clades.

\begin{figure}[t]
  \begin{center}
    \includegraphics[width=\linewidth]{\figpath/fig4_coherence_atteson.pdf}
  \end{center}
  \caption{\textbf{Where lineage recovery lives.} Mean true-clade recovery against
    \textbf{(A)} the fitted Atteson diagnostic $\widetilde m$ on a logarithmic x-axis and
    \textbf{(B)} label-free quartet
    additivity $A$. Large filled symbols summarize seven Flan and Llama simulations: full fine-tuning
    plus the five complete merged-weight LoRA conditions.
    Every mark is an average over the same 46 topology-eligible trees used in
    Table~\ref{tab:variants}; vertical intervals on simulation, PhyloLM, and
    subset summaries are $\pm1$ SE. Color identifies the seven primary training conditions only;
    auxiliary tensor subsets and individual $q/k/v$ matrices are shown in neutral gray so that their
    density does not create a second condition code. Dark dashed curves are pooled fractional-logit
    fits to all displayed estimator summaries; paired-bootstrap odds-ratio estimates accompany the
    reproducibility output.}
  \label{fig:coherence_recovery}
\end{figure}


\paragraph{Predicting held-out behavior.}
\label{sec:results:behavioral}

While estimating phylogeny from weights is useful for safety and provenance, it can be of even more use if similarity in weights indicates similarity in behavior. Inspired by \citet{cloud2025subliminal}'s subliminal-learning experiment, we test
whether models closer in weight space produce more similar outputs on held-out tasks, using translation
(on a tree that never trained on it) along with four probe tasks that, to our knowledge, neither model
was trained on (HellaSwag, ARC-Challenge, MMLU, TruthfulQA).
Table~\ref{tab:behavior_holdout} reports a significant negative random-effects correlation across all held-out
probes in Llama. Flan HellaSwag is instead significantly positive. Empty responses (EOS) are observed output and are represented as an
explicit empty-output semantic state in the primary analysis. In the Llama full-FT runs,
38.64\% of HellaSwag leaf outputs were empty, versus 0.02\% for TruthfulQA and 0.00\% for ARC-C
and MMLU. For HellaSwag, including immediate EOS gives $r_{\mathrm{DL}}=-0.17$ [$-0.24,-0.10$], whereas excluding empty completions strengthens the all-node estimate to $-0.27$ [$-0.34,-0.19$] and leaves the other probe estimates essentially unchanged (App.~\ref{app:empty-sensitivity}). These results show that shared fine-tuning ancestry is associated with behavioral similarity on held-out probes. The near-zero Flan held-out-translation cell is instead a degeneration diagnostic.
Classification-trained leaves collapse to their label vocabulary rather than translate, so similarity
reflects whether two leaves fail in the same way, not transmission of held-out translation behavior
(App.~\ref{app:translation_probe}).



\begin{table}[t]
\caption{Weight$\to$behavior bridge under full fine-tuning: random-effects pooled
  $r_{\mathrm{DL}}$ with 95\% CI. Immediate EOS is represented explicitly; $^\ddagger$degenerate
  Flan probe. Heterogeneity is in Appendix~\ref{app:regression}.}
\label{tab:behavior_holdout}
\centering
\small
\setlength{\tabcolsep}{5pt}
\begin{tabular}{lcccccc}
\toprule
$r_{\mathrm{DL}}$ [95\% CI] & \multicolumn{2}{c}{Translation} & HellaSwag & ARC-C & MMLU & TruthfulQA \\
\cmidrule(lr){2-3}
Model & trained & held-out & held-out & held-out & held-out & held-out \\
\midrule
Llama-3.2-1B & \shortstack{$-0.20$\\{\scriptsize$[-.26,-.13]$}} & \shortstack{$-0.23$\\{\scriptsize$[-.29,-.17]$}} & \shortstack{$-0.17$\\{\scriptsize$[-.24,-.10]$}} & \shortstack{$-0.20$\\{\scriptsize$[-.26,-.13]$}} & \shortstack{$-0.19$\\{\scriptsize$[-.26,-.12]$}} & \shortstack{$-0.22$\\{\scriptsize$[-.27,-.16]$}} \\
\addlinespace
Flan-T5-base & \shortstack{$-0.10$\\{\scriptsize$[-.16,-.04]$}} & \shortstack{$+0.04^{\ddagger}$\\{\scriptsize$[-.01,+.09]$}} & \shortstack{$+0.09$\\{\scriptsize$[+.05,+.14]$}} & \shortstack{$-0.14^{\ddagger}$\\{\scriptsize$[-.19,-.09]$}} & \shortstack{$-0.16^{\ddagger}$\\{\scriptsize$[-.20,-.11]$}} & \shortstack{$-0.17^{\ddagger}$\\{\scriptsize$[-.22,-.12]$}} \\
\bottomrule
\end{tabular}
\end{table}


\paragraph{Comparing to behavior-only lineage.}
\label{sec:results:blackbox}

When weights are unavailable or different architectures need to be compared, lineage must be inferred from behavior alone. We compare our weight-based
estimation to PhyloLM~\citep{yax2025phylolm}, scored against the
same documented lineage. On the real Mistral family (Fig.~\ref{fig:real-trees}), PhyloLM and weights
both recover all three forks. Their informative-quartet additivity differs ($A=0.84$ vs.\ $0.98$), indicating that the weight-based distance matrix has a stronger additive signal; because both reconstructions are exact, this does not imply a topological advantage on this family.

The Mistral forks are large RLHF or direct-preference-optimization (DPO) weight changes that move behavior substantially. Our controlled trees are a more precise test. In the LoRA and low-update regimes fine-tuning barely shifts behavior, yet weight distance can still estimate phylogenies. Running PhyloLM on the shared trees against the same ground truth, behavior-only recovery falls well below weight-based recovery. Under full fine-tuning PhyloLM recovers $52\%$ of clades versus $100\%$ for weights, under rank-8 LoRA it recovers $35\%$ versus $100\%$, and under rank-64 LoRA it recovers $24\%$ versus $100\%$. Behavioral similarity is a real but attenuated and probe-dependent readout of weight distance ($|r_{\mathrm{DL}}|=0.17$--$0.23$ across interpretable held-out Llama probes), which helps explain why output-based distances can preserve some pairwise signal yet lose enough resolution to miss topology after small updates. Overall, weight-based estimation matches behavior-based estimation on this real family and substantially outperforms it in the controlled experiments.

\begin{figure}[t]
  \begin{center}
    \includegraphics[width=\linewidth]{\figpath/hf_mistral_hero_3panel.pdf}
  \end{center}
  \caption{\textbf{Documented Mistral lineage:} truth (left), our weight-based estimate (middle),
    and PhyloLM's output-based estimate (right)~\citep{yax2025phylolm}. Both recover all three splits.}
  \label{fig:real-trees}
\end{figure}
% -----------------------------------------------------------------------
\section{Discussion}
\label{sec:discussion}
% -----------------------------------------------------------------------

We show that model family phylogeny can be estimated from weights alone, and that using a small tensor subset---about 3\% of Flan's scalar weights in the encoder-key subset---estimates training trees as well as using all analyzed tensors. We find that clade recovery works across encoder--decoder and decoder-only models and survives LoRA training at low rank ($r=8$). Across our controlled conditions, PAER ($\mathrm{FN}=0$) ranges from 85--100\%. The label-free additivity proxy $A$ is associated with increased clade recovery in our experiments, although the relationship depends on architecture and tensor subset. In addition to estimating phylogeny and direct lineage from weights, we show that models closer in weight space tend to have more similar held-out outputs. The effect depends on the probe and ranges from $r_{\mathrm{DL}}=-0.23$ to $-0.17$ across interpretable held-out Llama probes. Finally, weight- and behavior-based estimation recover the same documented Mistral family, while weight-based estimation is substantially stronger on several controlled trees and requires no prompts.


We note that our experiments and methods have limitations.  One limitation is that our methods for phylogeny estimation using weights currently only work when the underlying architecture is identical, i.e., the weight space is the same.  This means that Llama cannot directly be compared to Mistral, for example. Another issue is that our controlled training tree generation results from a fairly simple process that may not result in the most realistic training trees.  We attempt to mitigate that with our real-world example. Future directions include better relating models across different architectures such as with CKA and further exploring the mechanism by which lineage is inherited. In summary, we hope that our work illuminates the need to consider models holistically rather than only as a collection of benchmark scores, and shows how analyzing weight space can shed light on model behavior.


% -----------------------------------------------------------------------
% Author acknowledgments omitted for double-blind review; add at camera-ready.

\subsection*{AI use statement}
Generative AI assistants were used to write and debug software across the experimental
pipeline (data processing, weight-distance computation, tree reconstruction and scoring, statistical
analysis, figures, and tables) and to review and edit the manuscript, including checking consistency between numerical results and prose,
improving readability, searching and summarizing relevant literature, and formatting citations.  The research questions, experimental design,
hypotheses, conceptual framework, and scientific interpretations are the authors' own. The authors
reviewed and tested AI-assisted code and independently verified all results, figures, numerical claims,
citations, and manuscript text. The authors take responsibility for the final content of this work,
including all text, claims, and artifacts produced with the aid of generative AI.

% -----------------------------------------------------------------------
\bibliography{references}
\bibliographystyle{iclr2027_conference}


% -----------------------------------------------------------------------
\appendix
% -----------------------------------------------------------------------

\section{Experimental Details}
\label{app:details}

\subsection*{Tree Generation}
\label{app:tree-gen}

Training trees are generated via a Poisson branching process as follows:

\begin{enumerate}
    \item The pretrained root model draws $k \sim \mathrm{Poisson}(\lambda = 1.5)$ children, with a minimum of 1. We instantiate the same topology with either Flan-T5-base or Llama-3.2-1B.
    \item Each subsequent node draws $k \sim \mathrm{Poisson}(\lambda = 1.5)$ children. Nodes that draw $k = 0$ naturally become leaves; no separate pruning step is needed.
    \item Breadth First Search expansion continues until either a node budget ($n_\mathrm{nodes} = 14$ training nodes, not counting the root) is exhausted or a maximum depth ($n_\mathrm{layers} = 4$) is reached. The budget cap creates the truncation; any node alive when the budget runs out becomes a leaf.
    \item Each non-root node is assigned one of the 36 task/dataset pairs (9 summarization, 10 classification, 8 QA, and 9 translation). The 36 pairs are randomly permuted and sampled without replacement within a tree, so datasets do not repeat but task families can. The topology uses seeds beginning at \texttt{20260707}; assignment seeds are 1--50.
    \item Recovery excludes trees with fewer than four leaves or no nontrivial truth split, because their unrooted topology contains no clade that can be meaningfully recovered.
\end{enumerate}

This procedure yields trees with 4--10 observed leaves across the 50 accepted replications, with variable depth and branching structure. It produces 641 trained nodes in total (365 leaves and 276 internal nodes), in addition to the untrained roots. The ground-truth topology is fully known by construction, enabling exact RF evaluation.

\paragraph{Sample accounting.} Each controlled grid targets 50 generated training trees, while
recovery and direct ordering have different eligibility rules. The Flan comparison clones the
same 50 topology IDs across all five training regimes. The 46 truths with at least one nontrivial
unrooted split enter recovery, and the same 26 enter
ordering in every row; the other 24 trees have no cross-root-branch comparison, so ordering is
undefined rather than missing. The Llama rank-8, rank-64, and full-FT regimes reuse those same
50 topology IDs and task/data assignments; the same 46 enter recovery and the same 26 enter ordering.
The PhyloLM comparison produces a behavioral reconstruction for all 50 trees in each matched-topology Llama condition, but its
topology summaries use the common 46-tree subset after excluding four truths with no nontrivial
unrooted split (topology IDs 015, 020, 037, and 047); matched weight values in
Table~\ref{tab:variants} are computed on those same 46 trees.
The Flan per-tensor profile in Fig.~\ref{fig:layertrace_grid} uses the same 46 topology-eligible
full-fine-tuning trees; the Llama panel uses the matched 46-tree
full-fine-tuning cohort.

\subsection*{Training Configuration}
\paragraph{Flan-T5-base.} We pin \texttt{google/flan-t5-base} to revision \texttt{7bcac572ce56\allowbreak db69c1ea7c8a\allowbreak f255c5d7c9672fc2}. Each node trained for exactly 2,000 steps without early stopping, with seed 42, learning rate $3 \times 10^{-4}$, batch size 8, gradient accumulation 4 (effective batch 32), warmup 200 steps, weight decay 0.01, and bf16 precision. Inputs and targets are capped at 512 and 128 tokens, respectively. We used a frozen subsample with seed 42 and up to 10,000 training and 1,000 evaluation examples per dataset, holding the sampled data fixed across nodes assigned the same dataset.

\paragraph{Llama-3.2-1B.} We used the base model (not Instruct), pinned to revision \texttt{4e20de362430\allowbreak cd3b72f300e6\allowbreak b0f18e50e7166e08}, with completion-only causal loss. Each node trained for exactly 2,000 steps without early stopping, with seed 42, batch size 8, gradient accumulation 4 (effective batch 32), 200 warmup steps, weight decay 0.01, bf16 precision, right padding, and a 1,024-token prompt-plus-completion window. Full fine-tuning used learning rate $2 \times 10^{-5}$; QKV LoRA used $3 \times 10^{-4}$, dropout 0.05, and $(r,\alpha)=(8,16)$ or $(64,128)$.

\subsection*{Behavioral Probes and Analysis Population}
The behavioral analysis uses every trained checkpoint in each tree: 641 nodes across 50
trees, including the 276 trained internal nodes as observed ancestors and excluding the untrained root.
The 365 terminal leaves form a separately reported sensitivity population. Weight distances are direct
cosine distances between the two observed checkpoints in each pair; internal nodes are not reconstructed
from their descendants. Translation uses 200 frozen prompts with one deterministic greedy completion per
node (up to 128 new tokens for Flan and 64 for Llama). HellaSwag, ARC-Challenge, MMLU, and TruthfulQA
each use 100 frozen prompts and three sampled completions per prompt (seeds 42, 43, and 44; temperature
1, $p=1$; up to 64 new tokens). Dolly, used only in the EOS sensitivity, uses 30 prompts, three draws,
and up to 96 new tokens.

Outputs are embedded with \texttt{all-MiniLM-L6-v2}, pinned to revision
\texttt{1110a243fdf4706b3f48f1d95db1a4f5529b4d41}. We average prompt-paired cosine similarity (and
prompt-by-draw similarity for the three-draw probes) before correlating it with weight distance within
each tree. The inference unit is the tree. Within-tree Fisher-$z$ variances are estimated by deleting one
trained node at a time, after which the tree correlations are pooled by DerSimonian--Laird. Pair counts
are descriptive and are not treated as independent observations.

\section{Layer Variant Definitions}
\label{app:layers}
Two families of variants appear in the paper. \emph{Tensor subsets} (Table~\ref{tab:layer_subsets})
restrict the full-weight distance to a chosen set of tensors, testing which tensors carry the lineage
signal; they are defined on Flan-T5-base (12 encoder + 12 decoder blocks). \emph{LoRA variants}
(Table~\ref{tab:lora_variants}) instead change which modules are fine-tuned, all at rank $r=8$,
$\alpha=16$, dropout $0.05$. Their controlled-tree analysis merges the full
root-to-node adapter chain so distances are computed between standalone full-weight leaves. The
all-projections row is retained as an explicitly marked cumulative-update ablation: it compares the
cumulative root-to-node adapter displacement with the shared base omitted, whereas the other LoRA rows
use standalone merged full weights. It is therefore excluded from the merged-weight diagnostic fit.

\begin{table}[h]
\caption{Tensor subsets (matched-topology Flan-T5-base full FT; cosine; $n=46$). ``SA.k'' is the self-attention key projection;
``pre-attention LayerNorm'' is the \texttt{layer.0.layer\_norm} preceding self-attention. The table's
percentages count tensors in the 282-tensor trained set, not scalar parameters. In particular,
\texttt{enc\_sak} is 12/282 tensors (4.3\%) but 7,077,888/247,577,856 scalar weights (2.86\%).
\emph{Clade} and PAER are polytomy-aware recovery percentages with SEs over the same runs.
\texttt{low\_signal} is the predefined LayerNorm control subset.}
\label{tab:layer_subsets}
\centering
\scriptsize
\setlength{\tabcolsep}{2pt}
\begin{tabular}{llrrrr}
\toprule
Variant & Tensors included & Count & \% tensors & Clade (SE) & PAER (SE) \\
\midrule
\texttt{full}         & all trained tensors                                   & 282 & 100          & 100.0 (0.0) & 100.0 (0.0) \\
\texttt{high\_signal} & SA.k $+$ pre-attn LN (encoder $+$ decoder)            & 48  & $\sim$17     & 100.0 (0.0) & 100.0 (0.0) \\
\texttt{sak}          & SA.k only (encoder $+$ decoder)                       & 24  & $\sim$8.5    & 100.0 (0.0) & 100.0 (0.0) \\
\texttt{encoder\_h}   & encoder SA.k $+$ pre-attn LN                          & 24  & $\sim$8.5    & 100.0 (0.0) & 100.0 (0.0) \\
\texttt{enc\_sak}     & encoder SA.k only                                     & 12  & 4.3           & 100.0 (0.0) & 100.0 (0.0) \\
\texttt{low\_signal}  & post-attn/cross-attn/post-FFN LNs (control)           & 36 & $\sim$13 & 100.0 (0.0) & 100.0 (0.0) \\
\bottomrule
\end{tabular}
\end{table}

\begin{table}[h]
\caption{LoRA target-module variants in the fixed-2,000-step Flan-T5-base cohort
($r=8$, $\alpha=16$, dropout $0.05$). Modules $q,k,v,o$ are attention
query/key/value/output projections; \texttt{wi\_0}, \texttt{wi\_1}, and \texttt{wo} are the two gated
FFN inputs and the FFN output. Counts follow the modules actually resolved by PEFT (36 modules per
attention projection and 24 per FFN projection); percentages use the 247,577,856 scalar-weight analysis
denominator.}
\label{tab:lora_variants}
\centering
\begin{tabular}{ll r r}
\toprule
Variant & Target modules & Trainable params & \% scalars \\
\midrule
\texttt{lora\_k\_r8}          & $k$                                  & 442,368   & 0.18 \\
\texttt{lora\_qkv\_r8}        & $q,k,v$                              & 1,327,104 & 0.54 \\
\texttt{lora\_qkvo\_r8}       & $q,k,v,o$                            & 1,769,472 & 0.71 \\
\texttt{lora\_all\_proj\_r8}  & $q,k,v,o,\mathrm{wi}_{0,1},\mathrm{wo}$ & 3,391,488 & 1.37 \\
\bottomrule
\end{tabular}
\end{table}

\subsection*{Per-Tensor Recovery}
\label{app:layertrace}
Treating each tensor as an independent estimator, Fig.~\ref{fig:layertrace_grid} reports the per-tensor
FN-rate across architectures. On the Flan trees the signal is diffuse across individual tensors:
attention key matrices and the remaining tensors each miss about 4\% of true clades on average. This
does not conflict with Table~\ref{tab:layer_subsets}, where averaging only the 12 encoder
\textsc{SelfAttention.k} matrices matches the all-weight reconstruction. \textsc{LayerNorm} produces
prominent local maxima in both architectures.

\begin{figure}[h]
  \begin{center}
    \includegraphics[width=\linewidth]{\figpath/fig3_layertrace_grid.pdf}
  \end{center}
  \caption{Block-aligned per-tensor FN-rate (\% of true clades missed; mean $\pm$ 2\,SE,
    cosine, full fine-tuning) for the encoder and decoder of Flan-T5-base ($d{=}768$, $n=46$)
    and the decoder of the matched Llama-3.2-1B cohort ($d{=}2048$, $n=46$). Each transformer block
    occupies one x-position; gray points show the remaining tensors without connecting unrelated
    tensor families. Flan key matrices and other tensors each average about 4\% missed clades, whereas
    the Llama profile is more key-concentrated (1.1\% vs.\ 24.0\%). Color and shape highlight
    self-attention keys, cross-attention keys, and \textsc{LayerNorm}. The two Flan block-0
    relative-attention-bias tensors (32.8\% encoder; 55.4\% decoder) are omitted from the plot;
    panel y-axes are 0--5\%, 0--30\%, and 0--100\%, respectively.}
  \label{fig:layertrace_grid}
\end{figure}

\section{Blackbox Experiment Details}
\label{app:blackbox}

\paragraph{PhyloLM reproduction.} For the behavior-only comparison on the documented Mistral family
(Fig.~\ref{fig:real-trees}), we reproduce PhyloLM \citep{yax2025phylolm} faithfully using its released
\texttt{math} gene bank and its published Nei-similarity genetic distance (upstream commit
\texttt{8c70edf}). We sample $G=128$ genes and, for each, draw $N=32$ raw completions of $4$ tokens per
model (temperature $1$, no chat template, matching PhyloLM's design), take the $4$-character allele,
form the per-gene allele populations, compute the Nei similarity between every pair of the ten
checkpoints, convert to a distance ($-\log$ of the similarity), and build the tree by neighbor-joining.
The reconstruction is scored against the same documented truth tree and with the same polytomy-aware
scorer used for the whitebox families.

The controlled comparison uses the same locked PhyloLM contract ($G=128$, $N=32$, four-token
continuations, temperature $1$) with sampling seed \texttt{20260803} on the 50 matched-topology Llama
rank-8 LoRA, rank-64 LoRA, and full-FT conditions. For LoRA leaves, we reconstruct and query the
ordered cumulative root-to-leaf adapter chain; full-FT leaves are loaded directly. Each condition
contains 50 completed reconstructions and 365 leaf populations, and all generated alleles are
nonempty. The four trees whose truths contain no nontrivial unrooted split are valid behavioral
reconstructions but cannot define clade recovery, PAER, RF, or FN, leaving the same $n=46$ topology
sample in every row of Table~\ref{tab:variants}.

\section{Additional Organic Lineage Checks}
\label{app:organic-additional}

The paper evaluates three documented organic HuggingFace lineages in total. The ten-checkpoint
Mistral-7B family is the main-text example because it contains three documented nontrivial splits and
supports the matched PhyloLM comparison. Two smaller checks establish that its whitebox result is not
unique to Mistral. In a seven-checkpoint Flan-T5-base lineage, averaged-matrix NJ and the BHV Fr\'echet
mean recover the sole documented Flan-Alpaca descendant clade (FN $=0/1$); its complementary split has
98.2\% Fr\'echet support under cosine distance. In a ten-checkpoint Llama-2-7B lineage, the cosine
Fr\'echet estimate likewise recovers the sole documented seven-model chat-descendant clade (FN
$=0/1$), with 95.2\% support for the complementary base/Nous-Hermes/Vicuna split. These two references
are bushy and specify no ordering within the recovered clade, so we claim recovery of one documented
split in each, not complete resolution of their descendant histories.


\section{Behavioral Meta-Analysis Diagnostics}
\label{app:regression}

\paragraph{Pooling estimator.} Because each tree has its own weight-space scale, a raw pooled slope is
not meaningful across trees. We standardize within tree and estimate the variance of each Fisher-$z$
correlation with a delete-one-trained-node jackknife, then pool tree-level effects using a
DerSimonian--Laird random-effects meta-analysis. Pair counts are descriptive and are never treated as
independent observations.

\paragraph{Per-cell heterogeneity.} Table~\ref{tab:behavior_het} gives, for every cell of the main-text
bridge (Table~\ref{tab:behavior_holdout}, paired semantic DV), $r_{\mathrm{DL}}$ with its $95\%$ CI,
the Higgins $I^2$, and the number of trees $k$. Heterogeneity ranges from low to substantial ($I^2$ $0$--$36\%$).
All cells except Flan translation-withheld exclude zero. The Llama cells and three flagged Flan cells
are negative, Flan HellaSwag is significantly positive, and Flan translation-withheld is near zero.
Flagged Flan probes are output-degeneration diagnostics rather than clean bridge evidence.

\begin{table}[h]
\caption{Per-cell heterogeneity for the weight$\to$behavior bridge (Table~\ref{tab:behavior_holdout};
  paired semantic DV with immediate-EOS completions represented explicitly, full fine-tuning).
  DerSimonian--Laird random-effects $r_{\mathrm{DL}}$ with $95\%$ CI, Higgins $I^2$,
  and number of trees $k$. $^\ddagger$degenerate Flan probe (App.~\ref{app:translation_probe}).}
\label{tab:behavior_het}
\begin{center}
\small
\begin{tabular}{llrrr}
\toprule
Model & Probe & $r_{\mathrm{DL}}$ [95\% CI] & $I^2$ & $k$ \\
\midrule
Llama-3.2-1B & Translation (trained)              & $-0.20$ [$-0.26,-0.13$] & 13\% & 50 \\
Llama-3.2-1B & Translation (held-out)             & $-0.23$ [$-0.29,-0.17$] & 0\% & 50 \\
Llama-3.2-1B & HellaSwag (held-out)               & $-0.17$ [$-0.24,-0.10$] & 26\% & 50 \\
Llama-3.2-1B & ARC-C (held-out)                   & $-0.20$ [$-0.26,-0.13$] & 33\% & 50 \\
Llama-3.2-1B & MMLU (held-out)                    & $-0.19$ [$-0.26,-0.12$] & 36\% & 50 \\
Llama-3.2-1B & TruthfulQA (held-out)              & $-0.22$ [$-0.27,-0.16$] & 19\% & 50 \\
Flan-T5-base & Translation (trained)              & $-0.10$ [$-0.16,-0.04$] & 0\% & 50 \\
Flan-T5-base & Translation (withheld)$^\ddagger$ & $+0.04$ [$-0.01,+0.09$] & 0\% & 50 \\
Flan-T5-base & HellaSwag (held-out)               & $+0.09$ [$+0.05,+0.14$] & 0\% & 50 \\
Flan-T5-base & ARC-C (held-out)$^\ddagger$         & $-0.14$ [$-0.19,-0.09$] & 0\% & 50 \\
Flan-T5-base & MMLU (held-out)$^\ddagger$          & $-0.16$ [$-0.20,-0.11$] & 0\% & 50 \\
Flan-T5-base & TruthfulQA (held-out)$^\ddagger$    & $-0.17$ [$-0.22,-0.12$] & 3\% & 50 \\
\bottomrule
\end{tabular}
\end{center}
\end{table}


\section{Immediate-EOS Sensitivity and Depth}
\label{app:empty-sensitivity}

An immediate EOS is an observed model response rather than a missing observation. The primary semantic
analysis therefore represents an empty completion as an explicit semantic state. Table~\ref{tab:eos-sensitivity}
reports the complementary sensitivity analysis that removes pairs involving an empty completion. The choice
matters materially only for HellaSwag, where immediate EOS is common.

\begin{table}[h]
\caption{Llama-3.2-1B full-FT all-trained-node semantic bridge with immediate-EOS outputs included as an explicit
  semantic state (primary) or excluded (sensitivity). Values are $r_{\mathrm{DL}}$ [95\% CI]. Empty
  rates are pooled over all leaf outputs and all three draws.}
\label{tab:eos-sensitivity}
\centering
\small
\begin{tabular}{lrrr}
\toprule
Probe & EOS included & EOS excluded & Empty outputs \\
\midrule
ARC-Challenge & $-0.195$ [$-0.257,-0.131$] & $-0.195$ [$-0.257,-0.131$] & $0.00\%$ \\
Dolly open-ended & $-0.081$ [$-0.134,-0.028$] & $-0.083$ [$-0.136,-0.030$] & $0.16\%$ \\
HellaSwag & $-0.175$ [$-0.245,-0.103$] & $-0.265$ [$-0.338,-0.189$] & $38.64\%$ \\
MMLU & $-0.189$ [$-0.256,-0.121$] & $-0.189$ [$-0.256,-0.121$] & $0.00\%$ \\
TruthfulQA & $-0.218$ [$-0.275,-0.161$] & $-0.219$ [$-0.275,-0.161$] & $0.02\%$ \\
\bottomrule
\end{tabular}
\end{table}

Table~\ref{tab:eos-depth} reports the same empty-output rate by leaf depth. HellaSwag empties increase
with depth, from 16.72\% at depth 1 to 40.25\% at depth 4. Immediate termination is therefore itself
associated with accumulated fine-tuning depth; excluding these outputs nevertheless strengthens the
negative all-node semantic association, as shown in Table~\ref{tab:eos-sensitivity}.

\begin{table}[h]
\caption{Immediate-EOS output rate by leaf depth in the Llama-3.2-1B full-FT evaluation trees. Rates
  pool outputs across all leaves observed at each depth and all three draws.}
\label{tab:eos-depth}
\centering
\small
\begin{tabular}{lrrrr}
\toprule
Probe & Depth 1 & Depth 2 & Depth 3 & Depth 4 \\
\midrule
ARC-Challenge & 0.00\% & 0.00\% & 0.00\% & 0.00\% \\
Dolly open-ended & 1.45\% & 0.34\% & 0.10\% & 0.06\% \\
HellaSwag & 16.72\% & 30.86\% & 41.13\% & 40.25\% \\
MMLU & 0.00\% & 0.00\% & 0.00\% & 0.00\% \\
TruthfulQA & 0.10\% & 0.04\% & 0.01\% & $<0.01\%$ \\
\bottomrule
\end{tabular}
\end{table}


\section{Matched Translation Holdout}
\label{app:translation_probe}

The no-translation cohort is matched to the ordinary 50-tree cohort. It
preserves every topology, node order, and non-translation assignment, replacing only the 162 nodes
originally assigned translation data with unused classification, QA, or summarization datasets. Thus
479 of 641 node assignments are unchanged, no translation assignments remain, and the cohort differs
from the ordinary full-FT cohort only where translation was removed. We use the same 200-prompt
translation endpoint to test whether weight distance predicts behavioral similarity when translation
was never a training task.

The primary all-trained-node estimate is
$r_{\mathrm{DL}}=+0.04$ [$-0.01,+0.09$] across 50 trees, so its interval includes zero. We retain this
cell as a \emph{degeneration diagnostic}, not as evidence that the weight--behavior bridge reverses.
Pairwise semantic similarity can be driven by whether models fail on the held-out task in the same
task-specific way rather than by genealogical proximity. In the matched Llama-3.2-1B
full-fine-tuning cohort, the translation bridge is negative and
strengthens when translation is withheld (Table~\ref{tab:behavior_holdout}). The near-zero Flan-T5
result remains the diagnostic described here rather than a clean reversal of the bridge.


\section{Task-Metadata Baseline}
\label{app:taskbaseline}

To test whether weight-based recovery merely reflects task identity, we define a baseline that uses
\emph{only} task-family labels and no weights. In the controlled cohort, every edge receives one
task/dataset pair; datasets are unique within a tree, but the four task families can repeat. For each
leaf we read its root-to-leaf sequence of task types from the manifest (e.g.\ summarization $\to$
summarization $\to$ translation). The metadata distance between two leaves is the length of the two
task-type sequences not covered by their shared prefix, $d(i,j)=(\lvert s_i\rvert-p)+(\lvert
s_j\rvert-p)$ with $p$ the shared-prefix length, an ancestry analog in which leaves sharing a longer
initial training-task sequence are placed closer.

Many leaves share identical task paths, so this metadata distance contains extensive ties and any NJ
reconstruction depends on tie handling. We therefore use it only as a qualitative diagnostic and do
not report a quantitative task-path benchmark. In the within-task design below, task-family identity is
constant; the sequence distance can retain differences in leaf depth, but it contains neither dataset
identity nor branch-specific task-family variation.

\paragraph{Verifying the organic reference lineages.} For the real HuggingFace families
(\S\ref{sec:results:weighttree}), the documented reference tree is taken from the checkpoints' model cards and
then corroborated against each repository's \texttt{config.json}: the declared base model, the tokenizer
vocabulary size, and the hidden size / layer count. For the Mistral-7B-v0.1 family this confirms all ten
checkpoints share the architecture ($4096$ hidden, $32$ layers) and separates the two vocabulary groups
($32000$ vs.\ $32002$, the latter adding control tokens), consistent with the documented fork structure.
We treat the reference as a documented-and-verified lineage, not as unverified metadata.


\section{Within-Task Recovery}
\label{app:withintask}

The controlled trees in the main text assign a task/dataset pair to each edge, so a reconstruction
could in principle exploit repeated task-family labels (App.~\ref{app:taskbaseline}). The separate
within-task experiment fixes every node in a tree to the \emph{same} task family and
varies the dataset per node (Flan-T5-base, full fine-tuning; one of summarization, QA, classification, or
translation). All task-path elements are therefore identical labels. Paths of different lengths still
encode leaf depth under the sequence distance above, but task-family metadata contains no dataset
identity or branch-specific task variation; recovery beyond a depth-only cue must come from
within-task, dataset-level displacement in the weights.

Table~\ref{tab:withintask} reports this auxiliary cohort separately from the controlled
experiments (polytomy-aware recovery, $\geq 4$ leaves, cosine). Recovery persists
without task diversity, though it is weaker and more task-dependent than the multi-task setting
($92$--$100\%$ clade, Table~\ref{tab:variants}): classification, QA, and translation recover $72$--$87\%$
of clades, while summarization is the weakest at $59\%$.
Since the leaves are task-identical, this recovery shows that lineage signal can still be estimated even within similar task settings.

\begin{table}[h]
\caption{Auxiliary within-task phylogeny recovery (Flan-T5-base full fine-tuning, one task family per tree with
  datasets varying per node; cosine, $\geq 4$ leaves). Polytomy-aware, as in
  Table~\ref{tab:variants}; SEs over runs (binomial for PAER).}
\label{tab:withintask}
\centering
\small
\begin{tabular}{lrrrrr}
\toprule
Task (all nodes) & $n$ & Clade rec. (SE) & PAER (SE) & RF & FN \\
\midrule
Classification & 38 & \textbf{87\%} (4\%) & \textbf{79\%} (7\%) & 1.87 & 0.24 \\
QA             & 47 & \textbf{76\%} (6\%) & \textbf{72\%} (7\%) & 2.19 & 0.38 \\
Translation    & 47 & \textbf{72\%} (6\%) & \textbf{64\%} (7\%) & 2.23 & 0.40 \\
Summarization  & 47 & \textbf{59\%} (6\%) & \textbf{47\%} (7\%) & 2.70 & 0.64 \\
\bottomrule
\end{tabular}
\end{table}

\end{document}
